%%%%%%%%%%%%%%%%%%%%%%%%%%%%%%%%%%%%%%%%%%%%%%%%%%%%%%%%%%%%%%%%%%%%%%
% Lớp 10 | Mệnh đề. Tập hợp | Mệnh đề
% Dạng 0D1?1-1: Xác định mệnh đề, mệnh đề chứa biến
% Số câu: 56  (sinh tự động bởi tach_cau.py; từ nay thầy sửa trực tiếp file này)
%%%%%%%%%%%%%%%%%%%%%%%%%%%%%%%%%%%%%%%%%%%%%%%%%%%%%%%%%%%%%%%%%%%%%%

%%%%%%%%%%%%%%%% TRẮC NGHIỆM 4 PHƯƠNG ÁN %%%%%%%%%%%%%%%%

% nguồn: MaHoaNH_TracNghiem_Lop10
\begin{ex}%[0D1N1-1]%[ma:aca19a3b]
	Trong các câu sau, câu nào  là mệnh đề?
	\choice {Bạn cao bao nhiêu?}{Hôm nay là thứ mấy}{\True Việt  Nam giáp Trung Quốc}{Con học bài đi!}
	\loigiai{ 
		\lq\lq Việt Nam giáp Trung Quốc\rq\rq~ là câu khẳng định đúng nên là  mệnh đề.
	} 
\end{ex}

% nguồn: MaHoaNH_TracNghiem_Lop10
\begin{ex}%[0D1N1-1]%[ma:17ef0977]
	Cho mệnh đề chứa biến $P(x)\colon$\lq\lq$x^{2}+5x+4=0$\rq\rq.  Mệnh đề nào dưới đây đúng?
	\choice {\True $P(-1)=0$}{ $P(-5)=0$}{$P(2)=0$}{$P(1)=0$} 
	\loigiai{ 
		Có  $P(-1)=(-1)^{2}+5\cdot(-1)+4=0$.
	 }
\end{ex}

% nguồn: MaHoaNH_TracNghiem_Lop10
\begin{ex}%[0D1N1-1]%[ma:3bae6548]
	Trong các câu sau,  câu nào là mệnh đề?
	 \choice {Bạn học trường nào}{\True Số 10 là số chẵn}{Hoa hồng đẹp quá}{ Các bạn có làm được bài kiểm tra này không}
	\loigiai{ 
		\lq\lq Số 10 là số chẵn\rq\rq~  là câu khẳng định đúng nên là mệnh đề.
	} 
\end{ex}

% nguồn: MaHoaNH_TracNghiem_Lop10
\begin{ex}%[0D1N1-1]%[ma:795b39d7]
	Phát biểu nào sau  đây là mệnh đề?
	\choice {Trời mưa to quá}{Lan có thích uống trà sữa không}{Bão Yagi  lớn quá}{\True Hà Giang là một tỉnh miền núi phía bắc của Việt Nam}
	 \loigiai{
		\lq\lq Hà Giang là một tỉnh miền núi phía bắc của Việt Nam\rq\rq~ là một câu khẳng định có  thể xác định được tính đúng sai (là đúng), nên là một mệnh đề. Các phát biểu còn lại là câu cảm thán hoặc câu hỏi, không phải mệnh đề.
	} 
\end{ex}

% nguồn: MaHoaNH_TracNghiem_Lop10
\begin{ex}%[0D1N1-1]%[ma:8cd7bac1]
	Cho mệnh đề chứa biến $P(n)\colon $\lq\lq$ n+1 \text{ chia hết  cho }4$\rq\rq. Mệnh đề nào dưới đây đúng?
	\choice {$P(1)$}{$P(4)$}{$P(6)$}{\True  $P(7)$}
	\loigiai{ 
		Ta có $P(7)=7+1=8$ chia hết cho  $4$.
	} 
\end{ex}

% nguồn: MaHoaNH_TracNghiem_Lop10
\begin{ex}%[0D1N1-1]%[ma:9f9316c6]
	Phát biểu nào sau  đây không là mệnh đề?
	 \choice
		{Trái  đất quay quanh Mặt Trời}
	{Đảo Yến ở tỉnh Quảng Bình}{Brazin ở Châu Phi}{\True Đảo Phú Quốc  thuộc tỉnh nào của Việt Nam}
	 \loigiai{
		Phát biểu \lq\lq Đảo Phú Quốc thuộc tỉnh  nào của Việt Nam?\rq\rq~ là câu hỏi, không phải là câu khẳng định nên không là mệnh đề.
	} 
\end{ex}

% nguồn: MaHoaNH_TracNghiem_Lop10
\begin{ex}%[0D1N1-1]%[ma:310557af]
	Cho mệnh đề chứa biến $P(x)\colon$\lq\lq $x+15\le x^2$\rq\rq\, với $x$ là số thực. Mệnh đề nào sau  đây là đúng?
	\choice 
		 {$P(0)$}
		{$P(3)$} 
		{$P(4)$} 
		{\True  $P(5)$}
	 \loigiai{
		$P(5)\colon$ \lq\lq $5+15\le 5^2$\rq\rq. 
	 }
\end{ex}

% nguồn: MaHoaNH_TracNghiem_Lop10
\begin{ex}%[0D1N1-1]%[ma:881a677e]
	Với giá trị nào của  $x$ thì \lq\lq $x^2-1 = 0$, $x \in \mathbb{N}$\rq\rq \,là mệnh đề đúng?
	\choice 
		{\True $x = 1$} 
		{$x =  -1$}
		{$x =\pm 1$} 
		{$x =  0$}
	\loigiai{ 
		Với $x =  1$ thì $1^2-1 = 0$.
	} 
\end{ex}

% nguồn: MaHoaNH_TracNghiem_Lop10
\begin{ex}%[0D1N1-1]%[ma:8d705c69]
	Câu nào sau  đây \textbf{không} là mệnh đề?
	\choice 
		{\True $x >  2$}
		{$3  < 1$}
		{$4 - 5 = 1$} 
		{Tam  giác đều là tam giác có ba cạnh bằng nhau}
	\loigiai{ 
		Vì $x > 2$ là mệnh đề chứa biến, không phải mệnh đề. 
	} 
\end{ex}

% nguồn: MaHoaNH_TracNghiem_Lop10
\begin{ex}%[0D1N1-1]%[ma:f5a6ecd3]
	Các kí  hiệu nào sau đây dùng để viết đúng mệnh đề \lq\lq $7$ là một số tự nhiên\rq\rq?
	\choice 
		{$7\subset \mathbb{N}$} 
		 {\True $7\in \mathbb{N}$}
		 {$7 < \mathbb{N}$}
		{$7 \leq \mathbb{N}$} 
	 \loigiai{
		Ký hiệu viết đúng $7\in \mathbb{N}$. 
	} 
\end{ex}

% nguồn: MaHoaNH_TracNghiem_Lop10
\begin{ex}%[0D1N1-1]%[ma:cc1b4652]
	Phát biểu nào sau đây là một  mệnh đề?
	 \choice
		 {Thời tiết hôm nay lạnh quá!}
		{Đề thi môn  Văn quá hay!}
		{\True Quảng Ninh là một tỉnh của  Việt Nam}
		{Số $-3$ có phải là số tự nhiên  không?}
	\loigiai{ 
		Phát biểu \lq\lq Quảng Ninh là một tỉnh của Việt Nam\rq\rq là một  mệnh đề.
	 }
\end{ex}

% nguồn: MaHoaNH_TracNghiem_Lop10
\begin{ex}%[0D1N1-1]%[ma:0c0f2bb7]
	Mệnh đề nào sau đây đúng? 
	\choice 
		{\True Một số vừa chia hết cho $2$ vừa chia hết  cho $3$ thì nó chia hết cho $6$}
		{$\forall x\in \mathbb{R}:\sqrt{x^2}=x$} 
		{Phương trình $x^2-2=0$ có nghiệm hữu tỉ} 
		{Hình thoi có hai đường chéo bằng nhau} 
	 \loigiai{
		Vì $2$ và $3$ là hai  số nguyên tố cùng nhau nên \lq\lq Một số vừa chia hết cho $2$ vừa chia hết cho $3$ thì nó chia hết cho $6$\rq\rq.
	 }
\end{ex}

% nguồn: MaHoaNH_TracNghiem_Lop10
\begin{ex}%[0D1N1-1]%[ma:34d18d1e]
	Trong các câu sau, câu nào  \textbf{không} phải là mệnh đề?
	 \choice
		{Hà Nội là thủ đô  của Việt Nam}
		{$7$  là số nguyên tố}
		{\True Buồn ngủ quá!} 
		{$\sqrt{2}$ là  số vô tỉ}
	\loigiai{ 
		Câu không phải  mệnh đề là \lq\lq Buồn ngủ quá!\rq\rq.
	} 
\end{ex}

% nguồn: MaHoaNH_TracNghiem_Lop10
\begin{ex}%[0D1N1-1]%[ma:5460b8cc]
	Trong các câu sau, câu nào \textbf{không} là  mệnh đề chứa biến?
	\choice 
		{\True Số  $2$ không phải là số nguyên tố}
		{$4x^2-x-5=0$} 
		 {$5x-2y=0$}
		{$2m+1$  chia hết cho $3$}
	\loigiai{ 
		\lq\lq Số  $2$ không phải là số nguyên tố\rq\rq không là mệnh đề chứa biến.
	 }
\end{ex}

% nguồn: MaHoaNH_TracNghiem_Lop10
\begin{ex}%[0D1N1-1]%[ma:48678a84]
	Câu nào trong các  câu sau không phải là mệnh đề?
	\choice 
		 {\True $\pi$ có phải là một số vô tỷ không?}
		{$6-2=2$} 
		{$\sqrt{5}$ là một số hữu  tỷ}
		{$\dfrac{8}{4}=3$} 
	 \loigiai{
		$\pi$ có phải là  một số vô tỷ không? là câu nghi vấn nên không phải là mệnh đề.
	} 
\end{ex}

% nguồn: MaHoaNH_TracNghiem_Lop10
\begin{ex}%[0D1N1-1]%[ma:d0926b8c]
	Trong các câu sau có bao nhiêu câu  là mệnh đề\\
	(1): Số $15$ là  một số chẵn.\\
	(2): $2x^2-3x+1=3$.\\ 
	(3): Các em hãy cố gắng  làm bài thi cho tốt.\\
	(4): Nếu  $1< 3$ thì $14< 7$.
	 \choice
		 {\True $2$}
		 {$3$}
		 {$1$}
		 {$4$}
	 \loigiai{
		Mệnh đề  là câu (1) và (4).
	 }
\end{ex}

% nguồn: MaHoaNH_TracNghiem_Lop10
\begin{ex}%[0D1H1-1]%[ma:1ef8c9a7]
	Cho mệnh đề chứa biến $P(n)\colon \lq\lq n^2-1$  chia hết cho $4$\rq\rq\, với $n$ là số nguyên. Xét xem các mệnh đề $P(5)$ và $P(2)$ đúng hay sai?
	 \choice
		 {$P(5)$ đúng và $P(2)$ đúng}
		{$P(5)$ sai và $P(2)$  sai}
		{\True $P(5)$ đúng và  $P(2)$ sai}
		 {$P(5)$ sai và $P(2)$ đúng}
	\loigiai{ 
		$P(5)=24$ đúng do $24\,\vdots\, 4$ còn $P(2)=3$ sai do $3$ không chia hết cho  $4$.
	} 
\end{ex}

% nguồn: MaHoaNH_TracNghiem_Lop10
\begin{ex}%[0D1H1-1]%[ma:2eb8042f]
	Câu nào trong các câu sau  \textbf{không} phải là mệnh đề?
	 \choice
		{$3+2=7$} 
		{$\forall x \in \mathbb{R}$, $x^2+1 > 0$} 
		{$-2< 0$} 
		 {\True $4+x$}
	\loigiai{ 
		Ta thấy $4+x$ là mệnh đề chứa biến, không  phải là mệnh đề.
	 }
\end{ex}

% nguồn: MaHoaNH_TracNghiem_Lop10
\begin{ex}%[0D1H1-1]%[ma:9a2c9950]
	Trong các câu sau, có bao nhiêu câu  là mệnh đề?\\
	 a) Hãy làm việc chăm chỉ!\\
	 b) Số $12$ chia hết cho $2$.\\
	c)  Số $11$ là số nguyên tố.\\
	d) Số thực $x$  là số chẵn.
	\choice 
		{$1$} 
		 {$4$}
		{\True  $2$}
		 {$3$}
	 \loigiai{
		 Khẳng định: “Số $12$ chia hết cho $2$” là mệnh đề đúng.\\
		Khẳng định: “Số $11$ là số nguyên tố”  là mệnh đề đúng.\\
		Khẳng định “Số thực $x$ là số chẵn”  không phải là mệnh đề mà đây là mệnh đề chứa biến.\\
		Hãy làm việc chăm chỉ! Đây là câu cảm  thán nên không phải là mệnh đề.\\
	Vậy các khẳng định  trên có $2$ mệnh đề.}
\end{ex}

% nguồn: MaHoaNH_TracNghiem_Lop10
\begin{ex}%[0D1K1-1]%[ma:b348ab5a]
	Mệnh đề "$\exists x\in\mathbb{R}:  x^2-3x+2=0$" được mô tả bởi mệnh đề nào dưới đây?
	\choice{Mọi số thực $x$ đều là nghiệm của phương  trình $x^2-3x+2=0$}
		{\True Có ít nhất một số thực $x$ là nghiệm  của phương trình $x^2-3x+2=0$}
		 {Có duy nhất một số thực $x$ là nghiệm của phương trình $x^2-3x+2=0$}
		{Nếu $x$ là số thực thì  $x^2-3x+2=0$}
	\loigiai{Mệnh đề "$\exists x\in\mathbb{R}: x^2-3x+2=0$" được mô tả lại như sau: "Có ít nhất một số thực $x$ là nghiệm của phương trình  $x^2-3x+2=0$".}
\end{ex}

% nguồn: MaHoaNH_TracNghiem_Lop10
\begin{ex}%[0D1B1-1]%[ma:1f6f56fe]
	 Câu nào trong các câu sau không phải là mệnh đề?
	\choice 
		{\True ``$\pi $  có phải là một số vô tỷ không?''}
		{``$2+2=5$''} 
		{``$\sqrt{2}$ là  một số hữu tỷ''}
		 {``$\dfrac{4}{2}=2$''}
	\loigiai{``$\pi  $ có phải là một số vô tỷ không?'' không là mệnh đề.}
\end{ex}

% nguồn: MaHoaNH_TracNghiem_Lop10
\begin{ex}%[0D1B1-1]%[ma:d99bf976]
	Cho mệnh đề $P(x)$: \lq\lq $x^2<x$\rq\rq. Mệnh đề nào sau đây  là đúng
	 \choice
		{$P(2)$} 
		 {$P(1)$}
		{\True $P\left( \dfrac{1}{2}\right)  $}
		 {$P(0)$}
	\loigiai{ 
		 $P\left( \dfrac{1}{2}\right)$ là mệnh đề đúng.
	} 
\end{ex}

% nguồn: MaHoaNH_TracNghiem_Lop10
\begin{ex}%[0D1B1-1]%[ma:60300e8c]
	Trong các câu sau, câu nào \textbf{không} phải là mệnh đề? 
	 \def\dotEX{}
	\choice  {\True Paris có phải là thủ đô của nước Pháp không?}
		{Paris  là thủ đô của nước Pháp.}
		 {$\sqrt{3}$ là một số vô tỉ.}
		{Tam giác $ABC$ có một  góc tù.}
	\loigiai{ 
		Câu a) không thể khẳng định đúng hoặc sai nên nó không phải là mệnh  đề.
	 }
\end{ex}

% nguồn: MaHoaNH_TracNghiem_Lop10
\begin{ex}%[0D1B1-1]%[ma:7eeadab8]
	Trong các câu sau,  câu nào {\bf không} phải là mệnh đề?
	 \choice
		 {\True Chị ơi, mấy giờ rồi?}
		{Bắc  Kinh là thủ đô của Việt Nam}
		{$\sqrt{2}$  là số hữu tỉ}
		{Số  $4$ là số chẵn}
	 \loigiai{
		Một câu hỏi thì  không phải là mệnh đề.
	} 
\end{ex}

% nguồn: MaHoaNH_TracNghiem_Lop10
\begin{ex}%[0D1B1-1]%[ma:0229015a]
	Câu nào trong các câu sau không phải là  mệnh đề?
	 \choice
		{Mọi số tự  nhiên đều là số nguyên}
		{  Số $2017$ là số nguyên tố}
		{ Tổng các góc trong của một tam  giác bằng $90^\circ$}
		{ \True $x^2-3x+2>0$} 
	 \loigiai{
		Câu “$x^2-3x+2>0$” vừa có thể đúng,  vừa có thể sai phụ thuộc vào giá trị của $x$, do đó nó không phải là mệnh đề.
	 }
\end{ex}

% nguồn: MaHoaNH_TracNghiem_Lop10
\begin{ex}%[0D1B1-1]%[ma:54374a85]
	Trong các câu sau, câu nào \textbf{không phải} là mệnh đề? 
	 \choice
		 {$5+2=8 $}
		 {$-2\le 0$}
		 {$ 4-\sqrt{17}>0$}
		{\True  $5+x=2 $}
	\loigiai{ 
		 \lq\lq$5+x=2 $\rq\rq\ không phải là mệnh đề.
	} 
\end{ex}

% nguồn: MaHoaNH_TracNghiem_Lop10
\begin{ex}%[0D1B1-1]%[ma:171e3430]
	Cho mệnh  đề chứa biến $P\left(x\right)\colon$ ``$x^2-3x+2=0$''. Mệnh đề $P\left(x\right)$ \textbf{đúng} khi nào?
	 \choice
		 {$x=0$}
		 {\True $x=1$}
		{$x=-1$} 
		 {$x=-2$}
	 \loigiai{
		Ta có $ x^2-3x+2=0\Leftrightarrow x=1\vee x=2 $  suy ra $ P(1) $ là mệnh đề đúng.
	 }
\end{ex}

% nguồn: MaHoaNH_TracNghiem_Lop10
\begin{ex}%[0D1B1-1]%[ma:2e078104]
	Cho mệnh đề chứa biến  $P(x)\colon$ ``$x\in\mathbb{R}\colon \sqrt{x}\ge x$''. Mệnh đề nào sau đây là \textbf{sai}?
	 \choice
		 {$P(0)$}
		 {$P\left(\dfrac{9}{16}\right)$}
		{$P\left(\dfrac{1}{4}\right)$} 
		{\True $P(2)$} 
	 \loigiai{
		Với $ x = 2 $, ta có $ \sqrt{2}<2 $ do đó $ P(2) $ là mệnh đề  sai.
	 }
\end{ex}

% nguồn: MaHoaNH_TracNghiem_Lop10
\begin{ex}%[0D1B1-1]%[ma:c8c863dd]
	Với giá trị nào của $x$ để mệnh đề chứa biến $Q(x)\colon$ ``$2x^2-5x+2=0$''  là một mệnh đề đúng?
	 \choice
		 {\True $x=\dfrac{1}{2}$}
		 {$x=1$}
		{$x=3$} 
		{$x=5$} 
	\loigiai{ 
		Ta có $  Q\left(\dfrac{1}{2}\right)\colon 0=0$ là mệnh đề đúng nên $ x=\dfrac{1}{2} $ là một giá trị thỏa mãn.
	 }
\end{ex}

% nguồn: MaHoaNH_TracNghiem_Lop10
\begin{ex}%[0D1B1-1]%[ma:88386dac]
	Cho mệnh đề chứa biến $P(x): $\lq\lq $ x  \in \mathbb{R} : \sqrt{x} \geq x$ \rq\rq. Mệnh đề nào sau đây là sai?
	 \choice
		 {$P\left(\dfrac{9}{16}\right)$}
		{$P\left(\dfrac{1}{4}\right)$} 
		 {$P(0)$}
		{\True $P(2)$} 
	 \loigiai{
		Vì $P(\dfrac{9}{16})$ là \lq\lq $\sqrt{\dfrac{9}{16}}\geq \dfrac{9}{16}$\rq\rq, đó là mệnh  đề đúng.\\
		 Vì $P(\dfrac{1}{4})$ là \lq\lq $\sqrt{\dfrac{1}{4}}\geq \dfrac{1}{4}$\rq\rq, đó là mệnh đề đúng.\\
		Vì  $P(0)$ là \lq\lq $\sqrt{0}\geq 0$\rq\rq, đó là mệnh đề đúng.\\
		Vì $P(2)$ là \lq\lq $\sqrt{2}\geq 2$ \rq\rq, đó là mệnh đề  sai.\\
	 }
\end{ex}

% nguồn: MaHoaNH_TracNghiem_Lop10
\begin{ex}%[0D1B1-1]%[ma:a6108115]
	Với giá trị nào của $x$ để mệnh đề chứa biến $Q(x)$: \lq\lq $2x^2-5x+2=0$\rq\rq\, là một mệnh  đề đúng?
	\choice 
		{$x=1$} 
		{\True  $x=\dfrac{1}{2}$}
		{$x=3$} 
		 {$x=5$}
	\loigiai{ 
		Với $x=\dfrac{1}{2}$ ta  có $Q\left(\dfrac{1}{2}\right)=0$ mệnh đề đúng.\\
		Với $x=1$ ta có $Q(1) =-1 \neq  0$ mệnh đề sai.\\
		Với $x=3$ ta có $Q(3)=5 \neq 0$  mệnh đề sai.\\
		Với $x=5$ ta có $Q(5)=27 \neq  0$ mệnh đề sai.
	 }
\end{ex}

% nguồn: MaHoaNH_TracNghiem_Lop10
\begin{ex}%[0D1B1-1]%[ma:ff4c55af]
	 Trong các câu sau, có bao nhiêu câu là mệnh đề?\\
	a) Hãy  mở cửa ra!\\
	b) Số 20 chia hết  cho 8.\\
	c) Số 17 là một số  nguyên tố.\\
	d) Bạn có thích chơi bóng đá  không?
	\choice 
		 {$1$}
		{\True  $2$}
		{$3$} 
		 {$4$}
	\loigiai{ 
		Câu a), d) không là mệnh đề vì không thể  khẳng định là đúng hay sai.\\
		Câu b), c) là các mệnh đề vì đó là những khẳng định có tính đúng hoặc sai. 
	} 
\end{ex}

% nguồn: MaHoaNH_TracNghiem_Lop10
\begin{ex}%[0D1B1-1]%[ma:f3e0f5a5]
	Trong các câu  sau, có bao nhiêu câu là mệnh đề?\\
	- Hãy cố  gắng học thật tốt!\\
	- Số $20$ chia hết cho $6$.\\ 
	-  Số $5$ là số nguyên tố.\\
	- Số  $x$ là số chẵn
	\choice 
		{$4$} 
		{$3$} 
		{\True $2$} 
		 {$1$}
	 \loigiai{
		Có  hai mệnh đề là\\
		- Số $20$  chia hết cho $6$.\\
	- Số $5$ là số nguyên tố.} 
\end{ex}

% nguồn: MaHoaNH_TracNghiem_Lop10
\begin{ex}%[0D1B1-1]%[ma:a5eae881]
	Câu nào sau đây \textbf{\textit{không phải}} là mệnh đề? 
	\choice 
		{$3-2=1$} 
		 {$2<\sqrt{3}$}
		{\True $2+x=3$} 
		 {$x^2+1>0$}
	\loigiai{ 
		Ta có 
		\begin{itemize}
			\item $2+x=3$ là mệnh đề chứa biến,  không phải mệnh đề.
			\item $2<\sqrt{3}$ là mệnh đề. 
			\item  $3-2=1$ là mệnh đề.
			\item $x^2+1>0$ vừa là  mệnh đề, cũng vừa là mệnh đề chứa biến.
		\end{itemize}
	} 
\end{ex}

% nguồn: MaHoaNH_TracNghiem_Lop10
\begin{ex}%[0D1B1-1]%[ma:fd61ebc8]
	Cho mệnh đề chứa biến $P(x)$: \lq\lq$x\in \mathbb{R}:\, x^2-3x+2=0$\rq\rq.  Mệnh đề $P(x)$ \textbf{\textit{đúng}} khi nào?
	\choice 
		{$x=0$} 
		{\True $x=1$} 
		{$x=-1$} 
		{$x=-2$} 
	\loigiai{ 
		Ta  có
		\begin{itemize}
			\item Với $x=0$ ta có  $P(0)=2 \neq 0$ mệnh đề sai.
			 \item Với $x=1$ ta có $P(1)=0$ mệnh đề đúng.
			\item Với  $x=-1$ ta có $P(-1)=6\neq 0$ mệnh đề sai.
			 \item Với $x=-2$ ta có $P(-2)= 17 \neq 0$ mệnh đề sai.
		\end{itemize}
	 }
\end{ex}

% nguồn: MaHoaNH_TracNghiem_Lop10
\begin{ex}%[0D1B1-1]%[ma:96922961]
	Cho mệnh đề chứa biến $P(x)$: \lq\lq$x\in \mathbb{R}:\, \sqrt{x}\geq x$\rq\rq.  Mệnh đề nào sau đây là \textbf{\textit{sai}}?
	\choice 
		{$P(0)$} 
		{\True $P(2)$} 
		 {$P\left(\dfrac{9}{16}\right)$}
		{$P\left(\dfrac{1}{4}\right)$} 
	\loigiai{ 
		 Ta có
		\begin{itemize}
			\item $P(0)$: \lq\lq$0 \in \mathbb{R}:\sqrt{0} \geq 0$\rq\rq\,  mệnh đề đúng.
			\item $P(2)$: \lq\lq$2\in \mathbb{R}:\sqrt{2} \geq 2$\rq\rq\,  mệnh đề sai.
			 \item $P\left(\dfrac{9}{16}\right)$: \lq\lq$\dfrac{9}{16}\in \mathbb{R}: \sqrt{\dfrac{9}{16}} \geq \dfrac{9}{16}$\rq\rq\, mệnh đề đúng.
			\item $P\left(\dfrac{1}{4}\right)$: \lq\lq$\dfrac{1}{4}\in \mathbb{R}: \sqrt{\dfrac{1}{4}} \geq \dfrac{1}{4}$\rq\rq\, mệnh đề  đúng.
		\end{itemize}
	 }
\end{ex}

% nguồn: MaHoaNH_TracNghiem_Lop10
\begin{ex}%[0D1B1-1]%[ma:2c54d81b]
	Trong các câu sau,  có bao nhiêu câu là mệnh đề?\\
	a) Đăk Lăk là  một thành phố của Việt Nam.\\
	b) Sông Hương chảy ngang qua  thành phố Huế.\\
	c) Hãy trả lời câu  hỏi này!\\
	d) $5+19=24$.\\ 
	e)  $6+81=25$.
	 \choice
		 { $1$}
		 { $2$}
		 { $3$}
		{ \True  $4$}
	\loigiai{ 
		Câu c) không là mệnh đề vì không thể khẳng  định là đúng hay sai.\\
		Câu a), b), d),  e) là các mệnh đề vì đó là những khẳng định có tính đúng hoặc sai.
	} 
\end{ex}

% nguồn: MaHoaNH_TracNghiem_Lop10
\begin{ex}%[0D1B1-1]%[ma:313c279c]
	Trong các khẳng  định sau, có bao nhiêu khẳng định là mệnh đề?
	\begin{itemize}
		\item  “$2+4=7$ ”.
		\item Học, học nữa, học mãi. 
		\item Hình chữ nhật  có hai đường chéo cắt nhau tại trung điểm mỗi đường.
		\item Tam giác có hai đường cao bằng  nhau là tam giác cân.
	\end{itemize}
	 \choice
		 {$4$}
		{\True $3$} 
		{$2$} 
		{$1$} 
	\loigiai{ 
		Khẳng định "Học,  học nữa, học mãi" không phải là mệnh đề vì không xác định được nó đúng hay sai.\\
		Còn lại các khẳng  định khác đều là mệnh đề.
	 }
\end{ex}

% nguồn: MaHoaNH_TracNghiem_Lop10
\begin{ex}%[0D1B1-1]%[ma:dc4b86f2]
	Trong các câu dưới đây có bao nhiêu  câu là mệnh đề chứa biến ?
	\begin{enumerate}
		\item  \lq\lq $ x+y>1 $\rq\rq;
		\item \lq\lq$2-\sqrt{5}<0$\rq\rq; 
		\item  \lq\lq$4+x=3$\rq\rq;
		\item \lq\lq$\pi >3$\rq\rq. 
	\end{enumerate}
	\choice 
		 {$3$}
		 {$4$}
		 {\True $2$}
		 {$1$}
	\loigiai{ 
		Ta có (1) và (3) là các mệnh đề chứa biến còn (2) và (4) là mệnh đề.\\ 
		 Vậy có $2$ câu là mệnh đề chứa biến.
	 }
\end{ex}

% nguồn: MaHoaNH_TracNghiem_Lop10
\begin{ex}%[0D1B1-1]%[ma:637f2666]
	Cho các câu sau đây 
	\begin{enumerate}[(I)]
		\item “Phan-xi-păng là ngọn núi cao nhất Việt  Nam”.
		\item  “$\pi^2<9,86$ ”.
		 \item “Mệt quá!”.
		\item “Chị ơi,  mấy giờ rồi?”
	\end{enumerate}
	Hỏi có bao nhiêu câu  là mệnh đề?
	 \choice
		{$1$} 
		{$3$} 
		 {$4$}
		{\True  $2$}
	 \loigiai{
		Mệnh  đề là một khẳng định có tính đúng hoặc sai, không thể vừa đúng vừa sai.\\
	Do đó,  (I), (II) là mệnh đề, (III), (IV) không là mệnh đề.}
\end{ex}

% nguồn: MaHoaNH_TracNghiem_Lop10
\begin{ex}%[0D1B1-1]%[ma:f0e9857d]
	Trong các câu sau, có bao nhiêu câu là  mệnh đề?
	\begin{enumerate}
		 \item Hãy đi nhanh lên!
		\item Hà Nội  là thủ đô của Việt Nam.
		\item  $5+3=7$.
		\item $x+1=3$. 
		\item Năm $2020$ là  năm nhuận.
	\end{enumerate}
	\choice 
		{\True $3$} 
		{$2$} 
		{$4$} 
		 {$5$}
	 \loigiai{
		$\bullet$ Câu: Hãy đi nhanh lên! là câu cảm thán  không phải mệnh đề.\\
		$\bullet$ Câu: $x+1=3$ là mệnh đề chứa  biến không phải mệnh đề.
	} 
\end{ex}

% nguồn: MaHoaNH_TracNghiem_Lop10
\begin{ex}%[0D1B1-1]%[ma:e8ed8c20]
	Tìm số câu là mệnh đề trong  các câu sau.
	\begin{enumerate}[(I)]
		\item Ôi, mệt  quá!
		\item Số $ 3 $ là  số nguyên tố.
		\item Số $ 23  $ chia hết cho $ 3 $.
		\item Số  $x$ là một số lẻ.
	\end{enumerate}
	\choice 
		{\True $2$} 
		 {$1$}
		 {$3$}
		 {$4$}
	 \loigiai{
		(I) là  câu cảm thán nên không phải mệnh đề;\\
		(II) là  mệnh đề đúng;\\
		(III) là mệnh đề sai;\\ 
		(IV) là mệnh đề chứa biến.\\ 
		Tóm lại, chỉ có (II) và (III) là  mệnh đề.
	 }
\end{ex}

% nguồn: MaHoaNH_TracNghiem_Lop10
\begin{ex}%[0D1B1-1]%[ma:49bf32de]
	Trong  các câu sau, có bao nhiêu câu là mệnh đề?
	\begin{enumerate}
		\item Hãy mở cửa  ra!
		\item Số $20$ chia  hết cho $3$.
		\item Số $13$ là một số nguyên tố. 
		\item  Bạn có thích học Toán không?
	\end{enumerate}
	\choice 
		{  $1$}
		{ \True  $2$}
		 { $3$}
		{  $4$}
	 \loigiai{
		Ta  có
		\begin{enumerate}
			\item Hãy mở cửa ra! $\rightarrow$  Không phải mệnh đề.
			\item Số $20$ chia hết cho $3$. $\rightarrow$ Là mệnh đề. 
			\item Số $13$ là một số nguyên tố. $\rightarrow$  Là mệnh đề.
			\item Bạn  có thích học Toán không? $\rightarrow$ Không phải mệnh đề.
		\end{enumerate}
		Vậy có $2$  câu là mệnh đề.
	} 
\end{ex}

% nguồn: MaHoaNH_TracNghiem_Lop10
\begin{ex}%[0D1B1-1]%[ma:dc8cf601]
	Trong các câu sau, có bao nhiêu câu \textbf{\textit{không}} là mệnh đề? 
	\begin{enumerate}
		\item Hãy cố  gắng học thật tốt!
		\item Số $2020$  là số nguyên tố.
		\item Số $2019$ là số  lẻ.
		\item Mặt trời mọc hướng Tây. 
	\end{enumerate}
	 \choice
		 {\True $1$}
		{$2$} 
		 {$3$}
		{$4$} 
	 \loigiai{
		Ta có 
		\begin{enumerate}
			\item Hãy cố gắng học thật tốt! $\rightarrow$  Không là mệnh đề.
			\item  Số $2020$ là số nguyên tố. $\rightarrow$ Là mệnh đề.
			\item Số $2019$ là số  lẻ. $\rightarrow$ Là mệnh đề.
			\item Mặt trời mọc hướng Tây.  $\rightarrow$ Là mệnh đề.
		\end{enumerate}
		Vậy  có $1$ câu không là mệnh đề.
	 }
\end{ex}


%%%%%%%%%%%%%%%% ĐÚNG/SAI %%%%%%%%%%%%%%%%

% nguồn: MaHoaNH_DungSai_Lop10
\begin{ex}%[0D1N1-1]%[ma:c64b6828]
	Các câu sau  là đúng hay sai?
	\choiceTF 
		{\lq\lq$Q(x)$: $2x=1$\rq\rq\ không  là mệnh đề chứa biến}
		{\True \lq\lq$P(n)$:  $n$ chia hết cho $5$\rq\rq\ là mệnh đề chứa biến}
		{\True \lq\lq$S(x,y,z)$: $x^2+y^2=z^2$\rq\rq\ là mệnh đề chứa  biến}
		{\lq\lq$R(x,y)$: $2x+y=3$\rq\rq\ không  là mệnh đề chứa biến}
	\loigiai{ 
	\begin{itemchoice}
		\itemch Câu \lq\lq$Q(x)$: $2x=1$\rq\rq\  là mệnh đề chứa biến.
		\itemch Câu \lq\lq$P(n)$: $n$ chia hết cho $5$\rq\rq\ là  mệnh đề chứa biến.
		\itemch Câu \lq\lq$S(x,y,z)$: $x^2+y^2=z^2$\rq\rq\ là mệnh đề chứa biến. 
		\itemch Câu \lq\lq$R(x,y)$: $2x+y=3$\rq\rq\ là mệnh  đề chứa biến.
	\end{itemchoice}
	} 
\end{ex}

% nguồn: MaHoaNH_DungSai_Lop10
\begin{ex}%[0D1N1-1]%[ma:bd721c9d]
	Các câu sau là đúng hay sai? 
	 \choiceTF
		{\True \lq\lq$3$ là  số lẻ\rq\rq\ là mệnh đề}
		{\lq\lq$1+2>3$\rq\rq\  không là mệnh đề}
		{\lq\lq$\pi$ là số vô tỉ phải không?\rq\rq\ là mệnh đề} 
		{\True \lq\lq Đến năm $2050$, con người sẽ đặt chân lên Sao Hỏa\rq\rq\  là mệnh đề}
	\loigiai{ 
	\begin{itemchoice}
		\itemch Câu  \lq\lq$3$ là số lẻ\rq\rq\ là mệnh đề.
		 \itemch Câu \lq\lq$1+2>3$\rq\rq\ là mệnh đề.
		\itemch Câu \lq\lq$\pi$ là số vô tỉ phải không?\rq\rq\ không là mệnh đề. 
		\itemch Câu \lq\lq Đến năm $2050$, con người sẽ đặt chân lên Sao Hỏa\rq\rq\ là mệnh  đề.
	\end{itemchoice}
	} 
\end{ex}


%%%%%%%%%%%%%%%% TRẢ LỜI NGẮN %%%%%%%%%%%%%%%%

% nguồn: MaHoaNH_TraLoiNgan_Lop10
\begin{bt}%[0D1H1-1]%[ma:0925d2ad]
	Cho mệnh đề chứa biến $P(x)\colon $\lq\lq $ x \leq 6 $ \rq\rq. Hỏi tổng tất cả các số  tự nhiên $x$ làm cho mệnh đề $P(x)$ đúng là?
	 \shortans{21}
	 \loigiai{
		Ta có các số tự nhiên $ x $ sao cho mệnh đề $ P(x) $ đúng  là $ 0 $; $ 1 $; $ 2 $; $ 3 $; $ 4 $; $ 5 $; $ 6 $.\\
		Nên tổng cần tìm là $ 1+2+3+4+5+6=21 $. 
	 }
\end{bt}

% nguồn: MaHoaNH_TraLoiNgan_Lop10
\begin{bt}%[0D1H1-1]%[ma:fa012941]
	Xét mệnh đề chứa biến $P(x)\colon$ \lq\lq $x^3-3 x^2+2 x=0$ \rq\rq. Tính tổng bình phương tất cả các giá trị $ x $ thỏa mãn $ P(x) $ là mệnh đề  đúng.
	\shortans{5} 
	 \loigiai{
		Ta có $x^3-3 x^2+2 x=0 \Leftrightarrow\hoac{&x=0 \\ &x=1 \\  &x=2.}$
		Vậy tổng tất cả giá trị $ x  $ để mệnh đề $ P(x) $ đúng là $ 0^2+1^2+2^2=5 $.
	} 
\end{bt}


%%%%%%%%%%%%%%%% TỰ LUẬN %%%%%%%%%%%%%%%%

% nguồn: MaHoaNH_TuLuan_Lop10
\begin{bt}%[0D1B1-1]%[ma:c04ded8a]
	Trong các câu  dưới đây, câu nào là mệnh đề, câu nào là mệnh đề chứa biến?
	\begin{listEX}[1]
		\item Hãy  đi nhanh lên!
		 \item $5+7+4=15$.
		\item Năm $2002$ là năm nhuận. 
		\item Phương trình  $x^2-3x+2=0$ vô nghiệm.
		 \item $x^{10}-1$ chia hết cho $11$.
		\item Có vô số  số nguyên tố.
	\end{listEX}
	\loigiai{ 
		Các mệnh đề là 2, 3, 4,  5, 6.\\
		Các mệnh đề chứa biến là  5.
	} 
\end{bt}

% nguồn: MaHoaNH_TuLuan_Lop10
\begin{bt}%[0D1B1-1]%[ma:54d84bef]
	Trong các câu dưới đây, câu nào  là mệnh đề, câu nào là mệnh đề chứa biến?
	\begin{listEX}[1]
		\item Hà Nội là thủ đô của  Việt Nam.
		\item  Thượng Hải là một thành phố của Ấn Độ.
		 \item $1+x=2$.
		\item $27$ chia hết  cho $5$.
		\item  Hôm nay trời đẹp quá!
		\item Hôm nay là thứ mấy? 
	\end{listEX}
	 \loigiai{
		Các mệnh đề là 1, 2, 3, 4.\\ 
		Các mệnh  đề chứa biến là 3.
	} 
\end{bt}

% nguồn: MaHoaNH_TuLuan_Lop10
\begin{bt}%[0D1B1-1]%[ma:c3d2372f]
	Cho  các mệnh đề chứa biến:
	\begin{enumerate}[a)]
		\item $P(x)\colon$ \lq\lq$2x=1$\rq\rq; 
		\item $R(x,y)\colon$ \lq\lq$2x+y=3$\rq\rq\,  (mệnh đề này chứa hai biến $x$ và $y$);
		\item $T(n)\colon$ \lq\lq $2n+1$ là số chẵn\rq\rq\, ($n$ là số tự nhiên). 
	\end{enumerate}
	Với mỗi mệnh đề chứa biến  trên, tìm những giá trị của biến để nhận được một mệnh đề đúng và một mệnh đề sai.
	 \loigiai{
		\begin{enumerate}[a)]
			\item Với $x=\dfrac{1}{2}$ thì $P\left(\dfrac{1}{2}\right)\colon$ \lq\lq$2 \cdot \dfrac{1}{2}=1$\rq\rq\, là mệnh  đề đúng.\\
			Với $x=1$ thì $P(1)\colon$ \lq\lq$2\cdot1=1$\rq\rq\, là  mệnh đề sai.
			\item Với $x=1, y=1$ thì $R(1,1)\colon$  \lq\lq$2\cdot1+1=3$\rq\rq\, là mệnh đề đúng.\\
			Với $x=1, y=2$ thì $R(1,2)\colon$ \lq\lq$2\cdot1+2=3$\rq\rq\, là mệnh đề sai. 
			\item Lấy số tự nhiên $n_0$ bất kì ta đều được $2n_0+1$ là một số lẻ, nghĩa là $T\left(n_0\right)\colon$ \lq\lq$2 n_0+1$ là số chẵn\rq\rq\, là mệnh đề sai. Do đó, không có  giá trị $n_0$ của $n$ đề $T\left(n_0\right)$ là mệnh đề đúng. $T\left(n_0\right)$ là mệnh đề sai với số tự nhiên $n_0$ bất kì.
		\end{enumerate}
	 }
\end{bt}

% nguồn: MaHoaNH_TuLuan_Lop10
\begin{bt}%[0D1B1-1]%[ma:92d7156e]
	Tìm hai giá trị thực của $x$ để từ mỗi câu sau ta được một mệnh  đề đúng và một mệnh đề sai.
	\begin{multicols}{4}
		\begin{enumerate}[a.]
			 \item $x^2<x$.
			\item $x=5x$. 
			\item  $x^2>0$.
			\item $x>\dfrac{1}{x}$. 
		\end{enumerate}
	\end{multicols}
	 \loigiai{
		\begin{enumerate}[a.]
			\item Với  $x=\dfrac{1}{2}$ thì $\left(\dfrac{1}{2}\right)^2<\dfrac{1}{2}$ là mệnh đề đúng. Với $x=1$ thì $1^2<1$ là mệnh đề sai.
			\item Với $x=0$ thì $0=0\cdot 5$ là mệnh đề đúng. Với  $x=-1$ thì $-1=-1\cdot 5$ là mệnh đề sai.
			\item Với $x=3$ thì  $3^2>0$ là mệnh đề đúng. Với $x=0$ thì $0^2>0$ là mệnh đề sai.
			\item Với $x=2$ thì $2>\dfrac{1}{2}$ là mệnh đề đúng. Với $x=-1$ thì $-1>\dfrac{1}{-3}$ là mệnh đề sai. 
		\end{enumerate}
	 }
\end{bt}

% nguồn: MaHoaNH_TuLuan_Lop10
\begin{bt}%[0D1B1-1]%[ma:647b9cb3]
	Trong các câu dưới đây, câu nào là mệnh đề, câu nào là mệnh đề  chứa biến?
	\begin{enumEX}{1}
		\item Số $11$ là số  chẵn.
		\item Bạn có chăm học  không?
		\item Huế là một thành phố của Việt Nam. 
		\item  $2x+3$ là một số nguyên dương.
		\item $2-\sqrt{5}<0$. 
		 \item $4+x=3$.
		\item Hãy trả lời câu  hỏi này!
		\item Paris  là thủ đô nước Ý.
		\item Phương  trình $x^2-x+1=0$ có nghiệm.
		\item $13$ là một số nguyên  tố.
	\end{enumEX}
	\loigiai{ 
		Các mệnh đề là 1, 3, 4, 5,  6, 8, 9, 10.\\
		Các mệnh đề chứa biến là 4, 6. 
	} 
\end{bt}

% nguồn: MaHoaNH_TuLuan_Lop10
\begin{bt}%[0D1B1-1]%[ma:3cdef717]
	Trong các câu sau đây, câu  nào là mệnh đề?
	\begin{multicols}{2}
		\begin{enumerate}[a)]
			\item $3$  là số lẻ;
			\item $1+2>3$; 
			\item $\pi$ là số vô tỉ phải  không?
			\item $0{,}0001$ là số rất  bé;
			\item Đến  năm $2050$, con người sẽ đặt chân lên Sao Hoả.
		\end{enumerate}
	\end{multicols}
	 \loigiai{
		\begin{enumerate}[a)]
			\item \lq\lq$3$ là số lẻ\rq\rq\, là mệnh đề (mệnh đề đúng). 
			\item \lq\lq$1+2>3$\rq\rq\, là mệnh đề  (mệnh đề sai).
			\item \lq\lq$\pi$ là  một số vô tỉ phải không?\rq\rq\, là câu hỏi, không phải mệnh đề.
			\item  \lq\lq$0{,}0001$ là số rất bé\rq\rq\, không có tính hoặc đúng hoặc sai (do không đưa ra tiêu chí thế nào là số rất bé). Do đó, nó không phải là mệnh đề.
			\item \lq\lq Đến năm $2050$, con người sẽ đặt chân lên Sao Hoả\rq\rq\, là một khẳng  định chưa thể chắc chắn là đúng hay sai. Tuy nhiên, nó chắc chắn chỉ có thể hoặc đúng hoặc sai. Do đó, nó là một mệnh đề.
		\end{enumerate}
		\begin{note}
			Những mệnh đề liên quan đến toán học (như các mệnh đề ở câu \textbf{a)} và  \textbf{b)} trong \textbf{Ví dụ 1}) còn được gọi là mệnh đề toán học.
		\end{note}
	} 
\end{bt}

% nguồn: Ok_TuanLuan_Toan11_Lan1
\begin{bt}%[0D1Y1-1]%[ma:bbc4bfe7]
	Cho một tam giác đều $A B C$ cạnh $a$. Tam giác $A_1 B_1 C_1$ có các đỉnh là trung điểm các cạnh của tam giác $A B C$, tam giác $A_2 B_2 C_2$ có các đỉnh là trung điểm các cạnh của tam giác  $A_1 B_1 C_1, \ldots$, tam giác $A_{n+1} B_{n+1} C_{n+1}$ có các đỉnh là trung điểm các cạnh của tam giác $A_n B_n C_n, \ldots$ Gọi $p_1, p_2, \ldots, p_n, \ldots$ và $S_1, S_2, \ldots, S_n, \ldots$ theo thứ tự là chu vi và diện tích của các tam giác $A_1 B_1 C_1, A_2 B_2 C_2, \ldots, A_n B_n C_n, \ldots$.
	\begin{listEX}
		\item[a)] Tìm giới hạn của các  dãy số $\left(p_n\right)$ và $\left(S_n\right)$.
		\item[b)] Tìm các  tổng $p_1+p_2+\ldots+p_n+\ldots$ và $S_1+S_2+\ldots+S_n+\ldots$.
	\end{listEX}
	 \loigiai{
		\begin{listEX}
			\item[a)] Ta có $p_1, p_2, \ldots, p_n, \ldots$ lần lượt là chu vi của các tam giác $A_1 B_1  C_1, A_2 B_2 C_2, \ldots, A_n B_n C_n, \ldots$
			$$\begin{aligned}
			& p_1=3  a \\
			 & p_2=3 \cdot \frac{1}{2} a \\
			&  \ldots \\
			&  p_n=3 \cdot \frac{1}{2^{n-1}} a
		\end{aligned} $$
		suy ra $\lim p_n=\lim 3 \cdot \dfrac{1}{2^{n-1}} a=0$. 
		$$ \begin{aligned}
		& S_1=\frac{a^2 \sqrt{3}}{4}  \\
		&  S_2=\frac{1}{4} \frac{a^2 \sqrt{3}}{4} \\
		& \ldots  \\
		& S_n=\frac{1}{4^{n-1}} \cdot  \frac{a^2 \sqrt{3}}{4}
	\end{aligned}$$
	suy ra $\lim S_n=\lim  \dfrac{1}{4^{n-1}} \cdot \dfrac{a^2 \sqrt{3}}{4}=0$.
	\item[b)] Dựa vào dữ kiện đề bài suy ra tổng  $\left(p_n\right)$ là tổng của cấp số nhân lùi vô hạn với công bội $q=\dfrac{1}{2}$ và
	 $ p_1+p_2+\ldots+p_n+\ldots=\lim \left(p_n\right)=\dfrac{p_1}{1-q}=\dfrac{3 a}{1-\frac{1}{2}}=6a.$\\
	Dựa vào dữ kiện đề bài suy ra tổng $\left(S_n\right)$ là tổng của cấp số nhân lùi vô hạn với  công bội $q=\dfrac{1}{4}$ và
	$S_1+S_2+\ldots+S_n+\ldots=\lim \left(S_n\right)=\dfrac{S_1}{1-q}=\dfrac{\frac{a^2  \sqrt{3}}{4}}{1-\frac{1}{4}}=\dfrac{a^2 \sqrt{3}}{12}$.
	\end{listEX}
 }
\end{bt}

% nguồn: Ok_TuanLuan_Toan11_Lan2
\begin{ex}%[0D1Y1-1]%[ma:29be0bf8]
	Cho một tam giác đều $A B C$ cạnh $a$. Tam giác $A_1 B_1 C_1$ có các đỉnh là trung điểm các cạnh của tam giác $A B C$, tam giác $A_2 B_2 C_2$ có các đỉnh là trung  điểm các cạnh của tam giác $A_1 B_1 C_1, \ldots$, tam giác $A_{n+1} B_{n+1} C_{n+1}$ có các đỉnh là trung điểm các cạnh của tam giác $A_n B_n C_n, \ldots$ Gọi $p_1, p_2, \ldots, p_n, \ldots$ và $S_1, S_2, \ldots, S_n, \ldots$ theo thứ tự là chu vi và diện tích của các tam giác $A_1 B_1 C_1, A_2 B_2 C_2, \ldots, A_n B_n C_n, \ldots$.
	\begin{listEX}
		\item[a)] Tìm giới hạn của các  dãy số $\left(p_n\right)$ và $\left(S_n\right)$.
		 \item[b)] Tìm các tổng $p_1+p_2+\ldots+p_n+\ldots$ và $S_1+S_2+\ldots+S_n+\ldots$.
	\end{listEX}
	\loigiai{ 
		\begin{listEX}
			\item[a)] Ta có $p_1, p_2, \ldots,  p_n, \ldots$ lần lượt là chu vi của các tam giác $A_1 B_1 C_1, A_2 B_2 C_2, \ldots, A_n B_n C_n, \ldots$
			$$\begin{aligned}
			& p_1=3 a \\ 
			& p_2=3  \cdot \frac{1}{2} a \\
			& \ldots  \\
			&  p_n=3 \cdot \frac{1}{2^{n-1}} a
		\end{aligned} $$
		 suy ra $\lim p_n=\lim 3 \cdot \dfrac{1}{2^{n-1}} a=0$.
		$$ \begin{aligned}
		& S_1=\frac{a^2 \sqrt{3}}{4} \\ 
		& S_2=\frac{1}{4} \frac{a^2 \sqrt{3}}{4} \\ 
		&  \ldots \\
		& S_n=\frac{1}{4^{n-1}} \cdot \frac{a^2 \sqrt{3}}{4} 
	\end{aligned}$$
	suy ra $\lim S_n=\lim \dfrac{1}{4^{n-1}} \cdot \dfrac{a^2 \sqrt{3}}{4}=0$. 
	\item[b)] Dựa vào dữ kiện đề bài suy ra tổng $\left(p_n\right)$ là tổng của cấp số nhân lùi vô hạn với công bội $q=\dfrac{1}{2}$  và
	$ p_1+p_2+\ldots+p_n+\ldots=\lim \left(p_n\right)=\dfrac{p_1}{1-q}=\dfrac{3  a}{1-\frac{1}{2}}=6a.$\\
	Dựa vào dữ kiện đề bài suy ra tổng $\left(S_n\right)$ là tổng của  cấp số nhân lùi vô hạn với công bội $q=\dfrac{1}{4}$ và
	 $S_1+S_2+\ldots+S_n+\ldots=\lim \left(S_n\right)=\dfrac{S_1}{1-q}=\dfrac{\frac{a^2 \sqrt{3}}{4}}{1-\frac{1}{4}}=\dfrac{a^2 \sqrt{3}}{12}$.
	\end{listEX}
} 
\end{ex}
